\color{unidarkgreen}{\subsection{Логика деления стороны СН (ЛД СН)}}
\color{black}

\begin{enumerate}[label=\arabic{section}.\arabic{subsection}.\arabic*,  align=left, labelsep=4pt, leftmargin=0pt, itemindent=57pt]

\item
Ввод в работу логики деления средней стороны трансформатора от ступеней МТЗ СН осуществляется с
помощью уставок \textbf{«ЛД СН от МТЗ-1 СН»} и \textbf{«ЛД СН от МТЗ-2 СН»}.

\begin{figure}[H]
\centering
\includegraphics[width=0.88\textwidth,height=0.88\textheight,keepaspectratio]{\fbpath/022.1791 ЛД СН/img/LDSN.pdf}
\captionof{figure}{Функциональная схема ЛД стороны СН }\label{pict:LD_SN}
\end{figure}

\item
Оперативный ввод/вывод ЛД высокой стороны трансформатора производится с помощью БУ (блока
управления функциями) \textit{«Вывод ЛД СН»}. 

\item
Первым действием от направленных пусков ступеней МТЗ СН в сторону сети через выдержки времени
\textbf{«Тср МТЗ-1 откл СВ(ШСВ) СН»} и \textbf{«Тср МТЗ-2 откл СВ(ШСВ) СН»} логика деления формирует сигналы
предварительного деления с действием на отключение СВ (ШСВ) средней стороны трансформатора с
возможностью контроля включенного положения выключателя средней стороны.

\item
Вторым действием от направленных пусков ступеней МТЗ СН в сторону сети через выдержки времени
\textbf{«Тср МТЗ-1 откл В СН»} и \textbf{«Тср МТЗ-2 откл В СН»} логика деления формирует сигналы на отключение
выключателей средней стороны трансформатора.

\small
\setlength{\extrarowheight}{0.05cm} % высота ячеек в таблице
\begin{longtable}
{|>{\centering\arraybackslash}m{4.5cm}|>{\centering\arraybackslash}m{4.5cm}|>{\centering\arraybackslash}m{3cm}|>{\centering\arraybackslash}m{3.5cm}|}
\caption{Перечень уставочных параметров ЛД стороны СН\hfill\vspace{-0.5\baselineskip}}\label{ldsn:tbl1} \\ 
\hhline{|-|-|-|-|}
\rowcolor{gray!30}
Наименование уставки & Диапазон & Шаг &  Значение по умолчанию \\ 
\hhline{|-|-|-|-|}
\endfirsthead

\caption*{\hspace{3pt}\emph{Продолжение таблицы \ref{ldsn:tbl1}\hfill\vspace{-0.5\baselineskip}}} \\ % сделано по ГОСТ 2.105 п.6.8.7
\hhline{|-|-|-|-|}
\rowcolor{gray!30}
Наименование уставки & Диапазон & Шаг &  Значение по умолчанию \\  
\hhline{|-|-|-|-|}
\endhead
\hhline{|-|-|-|-|}

\endfoot
\endlastfoot

%%%%%%%%%%%%содержимое%%%%%%%%%%%%%%

\end{longtable}
\setlength{\extrarowheight}{0.0cm} % возврат параметра высота ячеек в таблице
\normalsize










































\end{enumerate} 